% Supplementary material for the Word-template manuscript. Built by: python paper/docx/build_docx.py --supp
% Same restricted LaTeX as content.tex; tables, figures and equations are numbered S1, S2, ...
% Every number comes from generated/*.tex. Do not type results by hand.

\title{Supplementary material}

\section{Paired household-level tests}
Table~\ref{tab:stats} lists every paired comparison. Each household's test MAE was averaged over the seeds of an arm
before testing, so each test has one pair per household (100 pairs). The Holm correction was applied over all rows of
the table together. Median $\Delta$MAE is the median over households of the difference A minus B; negative values
favour A.

\begin{table}
\caption{Wilcoxon signed-rank tests over households (seed-averaged MAE). A better: share of households with lower MAE
under A.}\label{tab:stats}
\begin{tabular}{llrrrr}
\toprule
A & B & median $\Delta$MAE (A$-$B) & A better (\%) & $p$ & $p_{\mathrm{Holm}}$ \\
\midrule
FedAvg, $\alpha$=0.3 & Local-only & +0.0013 & 34 & 0.0002 & 0.0023 \\
FedAvg, $\alpha$=0.3 & Centralised & +0.0026 & 0 & $<10^{-4}$ & $<10^{-4}$ \\
FedAvg, random & Centralised & +0.0021 & 0 & $<10^{-4}$ & $<10^{-4}$ \\
FedProx, $\alpha$=0.3 & FedAvg, $\alpha$=0.3 & -0.0001 & 75 & $<10^{-4}$ & $<10^{-4}$ \\
FedProx, cluster & FedAvg, cluster & +0.0001 & 36 & 0.0011 & 0.0103 \\
FedProx, $\alpha$=0.3, partial & FedAvg, $\alpha$=0.3, partial & +0.0003 & 24 & $<10^{-4}$ & $<10^{-4}$ \\
HA-FedProx, $\alpha$=0.3 & FedAvg, $\alpha$=0.3 & +0.0000 & 44 & 0.2480 & 0.4960 \\
HA-FedProx, $\alpha$=0.3 & FedProx, $\alpha$=0.3 & +0.0001 & 43 & 0.0031 & 0.0217 \\
HA-FedProx, cluster & FedAvg, cluster & +0.0004 & 27 & $<10^{-4}$ & $<10^{-4}$ \\
HA-FedProx, cluster & FedProx, cluster & +0.0001 & 39 & 0.0012 & 0.0103 \\
HA-FedProx, $\alpha$=0.3, partial & FedAvg, $\alpha$=0.3, partial & +0.0003 & 21 & $<10^{-4}$ & $<10^{-4}$ \\
HA-FedProx, $\alpha$=0.3, partial & FedProx, $\alpha$=0.3, partial & +0.0002 & 35 & 0.0469 & 0.1875 \\
FedAvg, cluster & FedAvg, random & +0.0001 & 45 & 0.0177 & 0.0884 \\
FedAvg, K=100 & Local-only & +0.0027 & 18 & $<10^{-4}$ & $<10^{-4}$ \\
FedAvg, K=100 & Centralised & +0.0045 & 0 & $<10^{-4}$ & $<10^{-4}$ \\
HA-FedProx, K=100 & FedAvg, K=100 & +0.0003 & 22 & $<10^{-4}$ & $<10^{-4}$ \\
FedAvg, $\alpha$=0.3, fp16 & FedAvg, $\alpha$=0.3 & -0.0002 & 92 & $<10^{-4}$ & $<10^{-4}$ \\
Centralised & Local-only & -0.0016 & 72 & $<10^{-4}$ & $<10^{-4}$ \\
FedAvg, K=1, random & Centralised & -0.0000 & 52 & 0.1046 & 0.3138 \\
FedAvg, random & FedAvg, K=1, random & +0.0023 & 1 & $<10^{-4}$ & $<10^{-4}$ \\
FedAvg, $\alpha$=0.3 + local FT & Local-only & -0.0004 & 52 & 0.3900 & 0.4960 \\
FedAvg, $\alpha$=0.3 + local FT & Centralised & +0.0010 & 27 & $<10^{-4}$ & $<10^{-4}$ \\
FedAvg, K=100 + local FT & Local-only & +0.0006 & 39 & 0.0142 & 0.0854 \\
\bottomrule
\end{tabular}

\end{table}

\section{Training budget and round selection}
Table~\ref{tab:budget} gives, for each number of clients, the local Adam steps per client and round, the round chosen
on validation error and the test MAE, for the 30-round budget of the main text and for runs of up to 100 rounds that
stopped after 25 rounds without a new best validation MAE. Fig.~\ref{fig:budget} shows the validation curves.

\begin{table}
\caption{Budget per number of clients $K$ (FedAvg, random grouping). Steps/round: Adam steps per client per round.
Best rd.: round selected on validation MAE (mean over seeds). Gap: test MAE relative to centralised.}\label{tab:budget}
\begin{tabular}{rrrrrlrr}
\toprule
$K$ & Steps/round & Best rd. (30) & MAE (30) & Gap (\%) & Best / run per seed (100) & MAE (100) & Gap (\%) \\
\midrule
1 & 4,535 & 18.6 & 0.0809 & −0.4 & -- & -- & -- \\
5 & 907 & 21.8 & 0.0826 & +1.8 & -- & -- & -- \\
20 & 227 & 22.0 & 0.0854 & +5.2 & 91, 98, 100 / 100, 100, 100 & 0.0832 & +2.5 \\
50 & 91 & 26.6 & 0.0870 & +7.2 & 52, 90, 25 / 77, 100, 50 & 0.0858 & +5.7 \\
100 & 46 & 29.2 & 0.0893 & +10.0 & 99, 95, 92 / 100, 100, 100 & 0.0871 & +7.3 \\
\bottomrule
\end{tabular}

\end{table}

\begin{figure}
\includegraphics[width=0.85]{revision_budget.png}
\caption{Validation MAE per round for each number of clients (mean over seeds). Solid: 30-round runs; dashed: runs of up
to 100 rounds. Dotted: centralised.}\label{fig:budget}
\end{figure}

\section{Heterogeneity-adaptive proximal coefficient}
The original submission proposed a rule that sets the proximal coefficient from the disagreement between client
updates. After round $r$ the server holds the updates $\Delta_k=w_k-w_r$ and their weighted mean $\bar{\Delta}$, and
computes
\begin{equation}\label{eq:rho}
\rho_r=\frac{\sum_k p_k\lVert\Delta_k-\bar{\Delta}\rVert^2}{\sum_k p_k\lVert\Delta_k\rVert^2}\in[0,1],\qquad \mu_{r+1}=\mathrm{clip}\left(\mu_0\rho_r,\ 10^{-4},\ 1\right).
\end{equation}
The rule never improved on FedAvg. Under the cluster partition it was worse for
\the\numexpr100-\StEonezerohafedproxstrongvsEsixfedavgstrongBetter\relax\,\% of households ($p_{\mathrm{Holm}}$
\pv{\StEonezerohafedproxstrongvsEsixfedavgstrongP}). Its signal $\rho_r$ ranked partitions by heterogeneity (Spearman
\RhoHSpearman{} with $H_{\mathrm{JS}}$, $p$ \pv{\RhoHSpearmanP}, $n=\RhoHRuns{}$ runs), but on average it rose from
\RhoRoundOne{} in round 1 to \RhoLast{} in the last round, a change about \RhoWithinOverBetween{} times its spread across
partitions. As the shared descent direction shrinks during training, any residual disagreement pushes $\rho_r$ towards
one, so a coefficient proportional to it grows late in training whatever the data. Table~\ref{tab:hetfull} gives all
three aggregation rules for each grouping.

\begin{table}
\caption{All aggregation rules for each grouping at $K=20$ (three seeds). $H$: composite index of the original
submission; $\bar\rho$: mean update disagreement.}\label{tab:hetfull}
\begin{tabular}{llrrrrr}
\toprule
Partition & Method & $H$ & $H_{JS}$ & $\bar\rho$ & Test MAE & Gap to central (\%) \\
\midrule
$\alpha$=0.1 & FedAvg & 0.242 & 0.0181 & 0.876 & 0.0856 $\pm$ 0.0005 & +5.6 \\
$\alpha$=0.1 & FedProx & 0.242 & 0.0181 & 0.863 & 0.0856 $\pm$ 0.0006 & +5.6 \\
$\alpha$=0.1 & HA-FedProx & 0.242 & 0.0181 & 0.844 & 0.0859 $\pm$ 0.0003 & +5.9 \\
$\alpha$=0.3 & FedAvg & 0.243 & 0.0153 & 0.874 & 0.0861 $\pm$ 0.0004 & +6.2 \\
$\alpha$=0.3 & FedProx & 0.243 & 0.0153 & 0.861 & 0.0858 $\pm$ 0.0001 & +5.9 \\
$\alpha$=0.3 & HA-FedProx & 0.243 & 0.0153 & 0.844 & 0.0863 $\pm$ 0.0003 & +6.5 \\
random & FedAvg & 0.255 & 0.0105 & 0.856 & 0.0853 $\pm$ 0.0009 & +5.2 \\
random & FedProx & 0.255 & 0.0105 & 0.845 & 0.0855 $\pm$ 0.0003 & +5.4 \\
random & HA-FedProx & 0.255 & 0.0105 & 0.822 & 0.0859 $\pm$ 0.0009 & +6.0 \\
$\alpha$=1.0 & FedAvg & 0.270 & 0.0124 & 0.867 & 0.0849 $\pm$ 0.0006 & +4.7 \\
$\alpha$=1.0 & FedProx & 0.270 & 0.0124 & 0.856 & 0.0857 $\pm$ 0.0001 & +5.7 \\
$\alpha$=1.0 & HA-FedProx & 0.270 & 0.0124 & 0.841 & 0.0857 $\pm$ 0.0008 & +5.8 \\
cluster & FedAvg & 0.381 & 0.0202 & 0.876 & 0.0857 $\pm$ 0.0003 & +5.7 \\
cluster & FedProx & 0.381 & 0.0202 & 0.866 & 0.0858 $\pm$ 0.0002 & +5.8 \\
cluster & HA-FedProx & 0.381 & 0.0202 & 0.851 & 0.0865 $\pm$ 0.0002 & +6.7 \\
\bottomrule
\end{tabular}

\end{table}

\section{Proximal coefficient grid}
Table~\ref{tab:tuning} gives the best validation MAE for each coefficient (cluster partition, seed 42). The row
$\mu=0$ is plain FedAvg.

\begin{table}
\caption{Best-round validation MAE (kWh) over the proximal coefficient grid.}\label{tab:tuning}
\begin{tabular}{lrr}
\toprule
$\mu$ (FedProx) or $\mu_0$ (adaptive) & FedProx & Adaptive (HA-FedProx) \\
\midrule
0 & 0.07238 & 0.07238 \\
0.0001 & 0.07260 & 0.07260 \\
0.001 & 0.07259 & 0.07268 \\
0.01 & 0.07320 & 0.07332 \\
0.1 & 0.07666 & 0.07682 \\
1 & -- & 0.07803 \\
\bottomrule
\end{tabular}

\end{table}

\section{Additional cohorts}
Two further cohorts of 100 homes were drawn with the same eligibility rules and stratification but different selection
seeds (43 and 44); they share three or four homes with the main cohort. Each arm ran with two seeds.

\begin{table}
\caption{Test MAE (kWh) of the main arms on three cohorts, and the FedAvg gap to centralised training at $K=1$, 5,
20 and 100.}\label{tab:cohorts}
\begin{tabular}{lrrrrrrl}
\toprule
Cohort & Centralised & Local-only & FedAvg, K=1 & FedAvg, K=5 & FedAvg, K=20 & FedAvg, K=100 & FedAvg gap, K=1/5/20/100 (\%) \\
\midrule
main (seed 42) & 0.0812 & 0.0838 & 0.0806 & 0.0828 & 0.0855 & 0.0898 & −0.7, +2.0, +5.4, +10.7 \\
cohort 43 & 0.0847 & 0.0878 & 0.0832 & 0.0853 & 0.0895 & 0.0909 & −1.7, +0.8, +5.7, +7.4 \\
cohort 44 & 0.0876 & 0.0866 & 0.0856 & 0.0866 & 0.0884 & 0.0910 & −2.3, −1.2, +0.9, +3.8 \\
\bottomrule
\end{tabular}

\end{table}

\section{Error by consumption quintile}
Households were divided into quintiles of mean training load (Q1 lowest). Table~\ref{tab:quint} gives the median
household test MAE in each quintile, Fig.~\ref{fig:equity} the distribution of the benefit of each arm over local
training, and Table~\ref{tab:equity} summary statistics of that benefit.

\begin{table}
\caption{Median household test MAE (kWh) by consumption quintile. FedAvg at $K=20$ unless stated.}\label{tab:quint}
\begin{tabular}{lrrrrr}
\toprule
Quintile & Local-only & Centralised & FedAvg, α=0.3 & FedAvg, α=0.3 + FT & FedAvg, K=100 \\
\midrule
Q1 & 0.0280 & 0.0285 & 0.0299 & 0.0297 & 0.0309 \\
Q2 & 0.0476 & 0.0474 & 0.0494 & 0.0481 & 0.0513 \\
Q3 & 0.0752 & 0.0741 & 0.0760 & 0.0782 & 0.0787 \\
Q4 & 0.1058 & 0.1038 & 0.1099 & 0.1073 & 0.1118 \\
Q5 & 0.1396 & 0.1319 & 0.1408 & 0.1393 & 0.1470 \\
\bottomrule
\end{tabular}

\end{table}

\begin{figure}
\includegraphics[width=0.7]{federation_benefit_cdf.png}
\caption{Distribution over households of the benefit of each collaborative arm over local training
($\mathrm{MAE}_{\mathrm{local}}-\mathrm{MAE}_{\mathrm{arm}}$, kWh; left of zero means worse off).}\label{fig:equity}
\end{figure}

\begin{table}
\caption{Equity relative to local-only training (seed-averaged household MAE): share harmed (any difference), median
benefit, 90th-percentile household MAE and share harmed after fine-tuning (FT).}\label{tab:equity}
\begin{tabular}{lrrrr}
\toprule
Arm & Harmed vs local (\%) & Median benefit (kWh) & P90 MAE & Harmed after FT (\%) \\
\midrule
Centralised LSTM (pooled raw data) & 28 & +0.0016 & 0.1364 & -- \\
FedAvg, random (IID-ish) & 56 & -0.0005 & 0.1439 & 46 \\
FedAvg, $\alpha$=0.3 & 66 & -0.0013 & 0.1445 & 48 \\
\bottomrule
\end{tabular}

\end{table}

\section{Calendar inputs and scaled error}
Sine and cosine encodings of time of day and day of week were added as input channels to every arm in
Table~\ref{tab:cal}, with all other settings unchanged.

\begin{table}
\caption{Effect of calendar inputs, and MASE (scaled per household by the in-sample one-step naive MAE). Gain: MAE
reduction relative to persistence; cal.: with calendar inputs. Gap: MAE relative to centralised training with the same
inputs.}\label{tab:cal}
\begin{tabular}{lrrrrrrr}
\toprule
Arm & MAE & MASE & Gain (\%) & MAE, cal. & MASE, cal. & Gain, cal. (\%) & Gap, cal. (\%) \\
\midrule
Centralised & 0.0811 & 0.987 & +8.9 & 0.0815 & 0.996 & +8.4 & -- \\
Local-only & 0.0840 & 1.006 & +5.6 & 0.0838 & 1.007 & +5.8 & -- \\
FedAvg, K=20 & 0.0854 & 1.045 & +4.0 & 0.0863 & 1.057 & +3.1 & +5.8 \\
FedAvg, K=100 & 0.0894 & 1.111 & −0.5 & 0.0919 & 1.158 & −3.3 & +12.8 \\
\bottomrule
\end{tabular}

\end{table}
